% !Mode:: "TeX:UTF-8"
\chapter{绪论}

\section{研究背景与意义}

随着集成电路设计规模的持续增长，寄存器传输级（RTL）代码的编写和验证已成为芯片设计流程中最耗时的环节之一。传统的RTL开发高度依赖工程师的专业经验——设计者需要根据功能规格手工编写Verilog或VHDL代码，并通过反复的编译、仿真和调试迭代来确保功能正确性。这一过程周期长、人力成本高，且容易因疏忽引入难以发现的逻辑错误。

近年来，大语言模型（Large Language Model, LLM）在软件代码生成领域展现出一定能力，能够根据自然语言描述生成多种编程语言的代码。这一进展自然引出了一个研究问题：大语言模型能否有效地生成硬件描述语言代码，从而辅助甚至部分替代人工的RTL开发过程？

然而，RTL代码生成与普通软件代码生成存在本质差异。RTL代码描述的是并行执行的硬件结构而非顺序执行的软件逻辑，正确的RTL设计不仅要求语法正确，还必须满足时序约束、位宽匹配、时钟域一致性等硬件特有要求，并通过EDA工具的编译和仿真验证。在零样本条件下，大语言模型生成的RTL代码往往存在编译错误、接口不匹配、边界条件遗漏等问题，尤其在涉及复杂状态机或浮点运算等场景时正确率较低。

这一现实使得将大语言模型的生成能力与EDA工具的验证能力相结合成为一个值得探索的方向。通过构建"生成—编译—仿真—反馈—修复"的迭代闭环，EDA工具的编译和仿真输出可以作为结构化的反馈信号引导模型定向修复代码错误，有望提升RTL代码的自动生成质量。

\section{国内外研究现状}

\subsection{LLM辅助RTL代码生成}

已有研究探索了使用大语言模型从自然语言规格直接生成Verilog代码的可行性\cite{chipgpt,thakur2023verigen}。在评估基准方面，NVlabs发布的VerilogEval\cite{verilogeval}提供了标准化的spec-to-RTL评估框架，香港科技大学发布的RTLLM\cite{rtllm}则提供了更高难度的设计任务集。这些基准的建立推动了RTL生成方法的标准化评估和横向比较。

现有研究表明，大语言模型在简单的组合逻辑和常见时序电路上具备一定的生成能力，但在涉及复杂算法、多状态FSM或精确位操作的设计上，单次生成的正确率仍有较大提升空间。

\subsection{EDA工具反馈与自动修复}

AutoChip\cite{autochip}等工作提出了利用EDA工具反馈驱动代码修复的方法，将传统的单次生成扩展为迭代修复过程。该类方法的核心思想是将编译器和仿真器输出的错误信息提取为结构化反馈，使模型能够进行定向修复而非盲目重新生成。RTLFixer\cite{rtlfixer}针对编译错误修复进行了专门探索，HDLDebugger\cite{hdldebugger}引入了检索增强技术辅助调试，这些研究共同体现了"验证反馈驱动修复"的研究趋势。

\subsection{现有工作的不足}

尽管上述研究在方法探索和概念验证方面取得了有价值的进展，仍有几个方面值得进一步分析。

在反馈收益的来源方面，反馈循环带来的通过率提升中有多少来自错误信息的引导、有多少仅来自多次重试的采样收益，现有工作较少对此进行严格的消融分析。在强模型与高难度场景方面，部分研究基于较早的模型或较简单的benchmark进行评估，需要新的实验证据。在反馈机制的影响因素方面，反馈粒度、反馈组织方式和初始提示词策略等维度对修复效果的影响，尚未得到系统的实验研究。此外，部分工作在实验结果管理和结论稳定性验证方面的机制不够完善。

\section{本文研究内容}

针对上述不足，本文的工作包含系统构建和实验研究两部分。

在系统构建方面，本文实现了一个AutoChip风格的RTL自动生成与反馈修复原型系统。该系统以Python编写，以Icarus Verilog作为EDA后端，实现了完整的``生成—编译—仿真—反馈—修复''闭环，并支持VerilogEval和RTLLM双benchmark适配、统一API网关下的多模型调用、五级反馈粒度和四种提示词策略的灵活组合。此外，对RTLLM 2.0全部50个设计进行了Icarus Verilog兼容性审计（41/50完全可用），从中精选12道高难度设计构成RTLLM\_STUDY\_12正式研究子集。

在实验研究方面，使用Haiku、Sonnet 4.6和GPT-5.4三个不同能力等级的模型，执行了七组递进式对照实验。实验从基础有效性验证起步，经由retry-only消融实验分离多采样与反馈信息各自的贡献，进而通过4轮重复实验检验结论的统计稳定性。在此基础上，分别探索了反馈粒度、多轮对话模式和提示词策略对修复效果的影响。全部实验均配有自动化产物管理和审计机制。

\section{本文主要贡献}

本文的贡献可从三个角度概括。

在工程实现方面，完成了一个可复现的AutoChip风格原型系统，支持VerilogEval与RTLLM双benchmark、统一API网关下的多模型调用、五级反馈粒度与四种提示词策略的组合配置。系统内置了实验产物自动序列化和数据审计脚本，确保每次实验可追溯至具体代码版本。围绕核心系统，项目还实现了RTLLM兼容性扫描、多模式实验运行和结果图表整理等配套脚本，使后续多组对照实验能够自动化执行和复现。

在实验验证方面，以RTLLM\_STUDY\_12作为核心任务集，在更高难度场景下验证了反馈循环的有效性。通过retry-only消融实验将反馈收益定量分解为多采样和错误反馈信息两部分，并通过4轮独立重复实验确认了结论的可复现性。实验中还观察到反馈效果的模型依赖现象：GPT-5.4上反馈始终带来正向增益，而Sonnet 4.6上反馈在当前配置下未能超越纯重试。

在机制探索方面，通过五级粒度实验揭示了反馈信息量与修复效率之间的倒U型关系——编译反馈和简洁反馈的稳健性最强，过于详尽的反馈可能引起信息过载而降低修复成功率。

\section{论文组织结构}

本文共分六章。第1章为绪论，介绍研究背景、现状和本文工作。第2章为相关技术与研究基础，介绍Verilog与RTL设计基础、大语言模型代码生成机制、EDA编译仿真流程和AutoChip方法。第3章为系统设计与实现，描述原型系统的架构和关键模块。第4章为实验设计，定义全部实验的方案和评价指标。第5章为实验结果与分析，呈现七组实验的结果并进行综合讨论。第6章为总结与展望，总结工作、讨论局限性并展望后续方向。
